\chapter{System Context}
\ChapterLabel{SystemContext}
\section{Chapter Responsibility}
\label{Section:SystemContextResponsibility}
Define the sensor-system boundary and its external relationships.

\noindent\textbf{Expected result:} A reviewable project record with explicit evidence and gaps.


\index{embedded software application}
\index{control app}
\index{sensor sample}
\section{Supported Boundary}
The system of interest contains the embedded software application and the host control app. The hardware sensor supplies sampled data to the embedded side; the control app receives that data on a host computer. The sensor and host hardware are environmental dependencies, not hardware developed in this documentation stage. Embedded software and control app are C++; the OS families named in the scope are macOS, Windows, and Linux. Their exact support matrix remains \hyperref[Question:Q-106]{Q-106}.

Figure~\ref{Figure:SensorContext} records only this supported data direction. Arrows do not select a protocol, connector, queue, thread model, or timing behavior. Bidirectional control has not been specified (\hyperref[Question:Q-105]{Q-105}).
\begin{figure}[htbp]
\centering
\begin{tikzpicture}[box/.style={draw=mbsBlue,fill=mbsLightBlue,align=center,text width=2.5cm,minimum height=1.5cm,font=\small},>=Latex]
\node[box] (sensor) at (0,0) {Hardware sensor\\selection open};
\node[box] (embedded) at (4,0) {Embedded software\\C++};
\node[box] (host) at (8,0) {Control app\\C++};
\node[draw=mbsGray,dashed,fit=(embedded)(host),inner sep=0.4cm,label={[font=\small]below:System of interest}] {};
\draw[->] (sensor.east) -- node[above=3pt,fill=white,inner sep=1pt,font=\scriptsize,text width=1.15cm,align=center]{sampled\\data} (embedded.west);
\draw[->] (embedded.east) -- node[above=3pt,fill=white,inner sep=1pt,font=\scriptsize,text width=1.15cm,align=center]{sampled\\data} (host.west);
\end{tikzpicture}

\caption{Sensor-to-control-app context; transport and hardware remain open.}
\label{Figure:SensorContext}
\end{figure}

\section{Stakeholders and Interfaces}
The supplied engineering owners are in \hyperref[Table:ProjectOwners]{the owner register}. The intended operator and mission need remain \hyperref[Question:Q-101]{Q-101}. Hardware, sample semantics, transport, constraints, failure behavior, and evidence remain \hyperref[Question:Q-102]{Q-102}, \hyperref[Question:Q-103]{Q-103}, \hyperref[Question:Q-104]{Q-104}, \hyperref[Question:Q-107]{Q-107}, \hyperref[Question:Q-108]{Q-108}. Course laboratory examples are context, not a selected sensor, FTDI device, or mandatory project architecture.
